% PDV Case Study — extracted from paper.tex for later use
% Original location: Results section, between top configs discussion and uncertainty section

\label{sec:pdv_casestudy}
% TODO: rewrite PDV case study paragraphs — reduce number regurgitation from table, remove "degrade gracefully", proper JoC style
To illustrate how deeply model and representation are entangled, we conducted a case study on Gaussian process (GP) models with two kernel choices---Tanimoto (a fingerprint kernel designed for binary features) and RBF (a general-purpose kernel for continuous features)---across two PDV representations: continuous (z-score normalised, 200 features) and binary (mean-threshold binarised, 200 features). In the main ANOVA, the GP was excluded because the Tanimoto kernel is fundamentally incompatible with continuous PDV: the kernel matrix collapses to near-zero for non-binary inputs, preventing model fitting. This exclusion is itself a manifestation of the interaction effect---the model's kernel constrains which representations are valid.

Table~\ref{tab:pdv_casestudy} summarises the results. On binary PDV, both kernels achieved comparable baselines ($R^2 \approx 0.874$--$0.875$), but their degradation profiles under noise differed markedly. The Tanimoto kernel degraded gracefully, maintaining $R^2 = 0.52$ at $\sigma = 0.5$ and $R^2 = 0.33$ at $\sigma = 0.9$, collapsing only at $\sigma = 1.0$. In contrast, the RBF kernel on the same binary data exhibited a sharp cliff: $R^2 = 0.58$ at $\sigma = 0.4$ followed by a collapse to $R^2 \approx 0$ at $\sigma = 0.5$. On continuous PDV, the RBF kernel achieved the strongest baseline of any model tested ($R^2 = 0.901$), matching XGBoost, but collapsed even earlier ($R^2 = 0.77$ at $\sigma = 0.1$, then $R^2 \approx 0$ at $\sigma = 0.2$).

This cliff-style failure is poorly captured by NDS, which fits a linear slope across all $\sigma$ levels. The Tanimoto GP's NDS ($-0.675$) and RBF GP's NDS on continuous data ($-0.688$) appear similar, but the underlying degradation profiles are qualitatively different: smooth versus catastrophic. This highlights a limitation of scalar robustness metrics and suggests that degradation curves should be inspected alongside summary statistics.

The case study also reveals that kernel choice matters more for GP robustness than representation choice. Tanimoto on binary PDV (NDS $= -0.675$) was more robust than RBF on continuous (NDS $= -0.688$) or binary (NDS $= -1.027$) PDV. However, this advantage comes at the cost of baseline performance: continuous PDV with an appropriate kernel yields $R^2 = 0.901$ versus $0.875$ for binary. This baseline--robustness tradeoff mirrors the broader pattern observed across architectures, but here it arises entirely from the kernel--representation interaction within a single model family.

\begin{table}[htbp]
\centering
\caption{GP kernel $\times$ PDV representation case study. Baseline $R^2$ and NDS (Gaussian noise, $\sigma = 0$--$1.0$) for two GP kernels on two PDV variants. $R^2$ at selected $\sigma$ levels shows degradation profiles. ``---'' indicates incompatible kernel--representation combination.}
\label{tab:pdv_casestudy}
\small
\begin{tabular}{llrrrrrr}
\toprule
\textbf{Kernel} & \textbf{PDV Variant} & \textbf{$R^2_0$} & \textbf{NDS} & \textbf{$R^2_{0.2}$} & \textbf{$R^2_{0.5}$} & \textbf{$R^2_{0.7}$} & \textbf{$R^2_{0.9}$} \\
\midrule
Tanimoto & binary (mean)  & $0.875$ & $-0.675$ & $0.681$ & $0.520$ & $0.432$ & $0.333$ \\
RBF      & binary (mean)  & $0.874$ & $-1.027$ & $0.680$ & $-0.005$ & $-0.004$ & $-0.004$ \\
RBF      & continuous     & $0.901$ & $-0.688$ & $-0.008$ & $-0.005$ & $-0.003$ & $-0.002$ \\
Tanimoto & continuous     & \multicolumn{6}{c}{--- (kernel incompatible with continuous features)} \\
\bottomrule
\end{tabular}
\end{table}

% Cross-references that pointed here:
% - Molecular Representations subsection: "For the GP kernel case study (Section~\ref{sec:pdv_casestudy})..."
% - Gauche GP description: "...we additionally tested a standard radial basis function (RBF) kernel within the same GP framework (Section~\ref{sec:pdv_casestudy})."
% - ANOVA methods: "The Gauche GP was also excluded because its Tanimoto kernel is incompatible with continuous-valued representations, preventing a fully crossed design (see Section~\ref{sec:pdv_casestudy})."
